\scanpage{7}
\mspage{\handwritten{- III -5-}}
\begingroup
\csname tagsleft@true\endcsname
le groupe des $p$-cochaînes \uline{alternées} de \leavevmode\rlap{C}\handwritten{\v{\phantom{C}}}ech, relativement au recouvrement $\underline U$, et à \deleted{\unsure{coefficients}} coefficients dans $\underline F$, il résulte de ce qui précède que l'on peut écrire
\[
C^p(\underline U,\underline F)=\mathop{\underset{\handwritten{\longrightarrow}}{\lim}}\limits_n C_n^p(M)\ .\tag*{(11)}
\]
Or, avec les notations du n\textsuperscript{o}1, $C_n^p(M)$ s'identifie à $K^{p+1}(\underline f,M)$, et l'application $\handwritten{\varphi}_{\unsure{nm|mn}}$ s'identifie à l'application $\handwritten{\varphi}_{\underline f^{n-m}}$ définie au n\textsuperscript{o}1. On a donc, pour tout $p\handwritten{\geq}0$, un isomorphisme canonique \deleted{\unsure{fonction}} fonctoriel en $M$
\[
C^p(\underline U,\underline F)\handwritten{\xrightarrow{\sim}}C^{p+1}((\underline f),M)\ .\tag*{(12)}
\]
La formule (3) du n\textsuperscript{o}1 montre en outre que ces isomorphismes sont compatibles avec les opérateurs cobord. Il en résulte ;

\uline{Proposition 2 .- Si} $X$ \uline{est} \inserted{\uline{un préschéma}} \uline{noethérien, ou} \deleted{\illegible{2}} \uline{un schéma don\unsure{t} l'espace de base est quasi-compact, il existe pour tout} $p\handwritten{\geq}1$ \uline{un isomorphisme canonique fonctoriel en} $\deleted{\unsure{M}}\handwritten{\underline F}$
\[
H^p(\underline U,\underline F)\handwritten{\xrightarrow{\sim}}H^{p+1}((\underline f),M)\ .\tag*{(13)}
\]
\uline{On a en outre une suite exacte fonctorielle en} $\deleted{\unsure{M}}\handwritten{\underline F}$
\[
0\handwritten{\to}H^0((\underline f),M)\handwritten{\to}M\handwritten{\to}H^0(\underline U,\underline F)\handwritten{\to}H^1((\underline f),M)\handwritten{\to}0\ .\tag*{(14)}
\]
Les relations (13) sont en effet évidentes d'après ce qui précède. D'autre part, \deleted{\unsure{on a} ${\unsure{H|C}}^{\unsure{0}}((\unsure{f}),\unsure{M})\unsure{=|,}{\unsure{H|C}}^{\unsure{1}}((\unsure{f}),\unsure{M})$} $M=C^0((\underline f),M)$ et \deleted{\illegible{1}} $C^0(\underline U,\underline F)={}=C^1((\underline f),M)$, donc $H^0(\underline U,\underline F)$ s'identifie au sous-groupe des 1-cocycles de $C^1((\underline f),M)$ ; la suite exacte (3) n'est autre alors que \inserted{celle qui résulte de} la définition des groupes de cohomologie $H^0((\underline f),M)$ et $H^1((\underline f),M)$.

\uline{Corollaire .- Supposons qu'il existe des} $g_i\handwritten{\in\Gamma}(U,O_X)$ \uline{tels que} $\handwritten{\sum_{i=1}^{\unsure{r|n}}}g_i f_{\unsure{i}}=1$ \inserted{\uline{au-dessus de} $U$}. \uline{Alors, pour tout faisceau quasi-cohérent} \deleted{\illegible{1}} $\deleted{\unsure{F}}\handwritten{\underline G}$ \uline{sur} $U$, \uline{on a} $H^i(\underline U,\underline G)=0$ \uline{pour} $i\handwritten{>}0$ ; \uline{si en outre} $X=U$, \uline{l'homomorphisme canonique} $M\handwritten{\to}H^0(\deleted{\unsure{M}}\underline{\unsure{U|M}},\underline F)$ \uline{est bijectif}. (les $X_{f_i}=U_i$ étant par hypothèse quasi-com\inserted{pacts)}

On peut en effet supposer $X=U$, et alors $H^i((\underline f),M)=0$ pour tout $i\handwritten{\geq\unsure{0|\alpha}}$ (n\textsuperscript{o}1, cor. de la prop.1); l'assertion découle alors aussitôt de (13) et (14). On observera que \inserted{puisque} $H^0(\underline U,\underline F)=H^{\deleted{\unsure{1}} \handwritten{0}}(U,\underline F)$, on redémontre ainsi le
\endgroup
% archive sheet number: 6
